% ----------------------------------------------------------------------------
%                                   CAPA
% ----------------------------------------------------------------------------
% Caso algum dos campos não se aplique ao seu trabalho, como por exemplo,
% se não houve coorientador, apenas deixe vazio.
% Exemplos:
% \coorientador{}
% \departamento{}

% ----------------------------------------------------------------------------
%                               DADOS DO TRABALHO
% ----------------------------------------------------------------------------
\titulo{Título do Trabalho} 
\subtitulo{\textnormal{subtítulo do trabalho}}
\autor{NOME COMPLETO D@ AUTOR(@)}
\autorcitacao{SOBRENOME, Nome} % Sobrenome em maiúsculo
\local{Brasília}
\data{2024}
\projeto{Trabalho de Conclusão de Curso}
\tituloAcademico{[bacharel ou licenciado ou tecnólogo]}

% ----------------------------------------------------------------------------
%                               DADOS DA INSTITUIÇÃO
% ----------------------------------------------------------------------------
\instituicao{Instituto Federal de Brasília}
\departamento{\textit{Campus} Brasília}
\programa{Tecnologia em Sistemas para Internet}
\logoinstituicao{0.15}{dados/figuras/ifb-logo.png}

% ----------------------------------------------------------------------------
%                              DADOS DOS ORIENTADORES
% ----------------------------------------------------------------------------
\orientador[Orientador:]{Nome do orientador} 
\instOrientador{[Departamento / Instituição a qual pertence]}
\tipoorientador{Presidente(a) / Orientador(a)}

\coorientador[Coorientadora:]{Nome da coorientadora}
\instCoorientador{[Departamento / Instituição a qual pertence]}

% Quando existir mais de um coorientador
% \coorientador[Coorientadores:]{Prof. Me. XXXXXX \newline
% Universidade XXXXXX - Câmpus XXXXXX \newline
% \newline Prof. Dr. XXXXXX. 
% \newline Universidade XXXXXX - Câmpus XXXXXX}
